% overleaf-compiler のサンプル原稿（LuaLaTeX）。
% Overleaf の「Download as source (.zip)」で落とした原稿と同じ形にしてある。
% 章は sections/、図は figures/、文献は references.bib に分けている。
\documentclass[a4paper,11pt]{ltjsarticle}

\usepackage{amsmath,amssymb}
\usepackage{graphicx}
\usepackage{booktabs}
\usepackage{tikz}
\usetikzlibrary{arrows.meta,positioning}
\usepackage[hidelinks]{hyperref}

\title{overleaf-compiler のサンプル原稿}
\author{overleaf-compiler}
\date{}

\begin{document}
\maketitle

\begin{abstract}
  この原稿は，overleaf-compiler の使い方を試すためのサンプルである．
  章ごとのファイル分割，図，表，数式，相互参照，文献の引用といった，
  Overleaf の原稿でよく使う要素を一通り含めている．
\end{abstract}

\section{はじめに}\label{sec:intro}

本稿では，ローカルの PC で LaTeX の原稿を直し，その場で PDF を確かめる流れを示す．
LaTeX の基本的な使い方は文献\cite{lamport1994latex,knuth1984texbook}に詳しい．
パッケージの使い方は文献\cite{mittelbach2004companion}にまとまっている．

本稿の構成は次のとおりである．\ref{sec:method}節で手順を，
\ref{sec:results}節で結果を示し，\ref{sec:conclusion}節でまとめる．

\section{手順}\label{sec:method}

手順の概要を図\ref{fig:flow}に示す．原稿を保存すると自動で組み直され，
PDF をダブルクリックするとソースの該当行へ移る．

\begin{figure}[t]
  \centering
  % 図1：原稿を直して PDF を確かめる流れ（TikZ）
\begin{tikzpicture}[
    box/.style={draw, rounded corners, minimum width=24mm, minimum height=9mm, align=center},
    arrow/.style={-{Stealth[length=2.5mm]}, thick}]
  \node[box] (edit) {原稿を直す};
  \node[box, right=12mm of edit] (build) {自動で組む};
  \node[box, right=12mm of build] (check) {PDF を確かめる};
  \draw[arrow] (edit) -- (build);
  \draw[arrow] (build) -- (check);
  \draw[arrow] (check.south) -- ++(0,-6mm) -| (edit.south);
\end{tikzpicture}

  \caption{原稿を直して PDF を確かめる流れ}
  \label{fig:flow}
\end{figure}

例として，$n$ 個の値 $x_1, \dots, x_n$ の平均 $\bar{x}$ と分散 $s^2$ を式\eqref{eq:mean}と式\eqref{eq:var}で定める．
\begin{align}
  \bar{x} &= \frac{1}{n} \sum_{i=1}^{n} x_i, \label{eq:mean} \\
  s^2 &= \frac{1}{n} \sum_{i=1}^{n} \left( x_i - \bar{x} \right)^2. \label{eq:var}
\end{align}

\section{結果}\label{sec:results}

表\ref{tab:values}に，式\eqref{eq:mean}と式\eqref{eq:var}で求めた値の例を示す．
数値は説明のために用意したものである．

\begin{table}[t]
  \centering
  \caption{平均と分散の例}
  \label{tab:values}
  \begin{tabular}{lrr}
    \toprule
    データ & 平均 $\bar{x}$ & 分散 $s^2$ \\
    \midrule
    A & 2.50 & 1.25 \\
    B & 3.00 & 2.00 \\
    C & 4.20 & 0.56 \\
    \bottomrule
  \end{tabular}
\end{table}

\section{おわりに}\label{sec:conclusion}

本稿では，overleaf-compiler で原稿を直して確かめる流れを示した．
直し終えた原稿は，一覧の項目のメニューにある「ダウンロード」で，Overleaf に戻せる zip にできる．


\bibliographystyle{plain}
\bibliography{references}

\end{document}
